\title{Thesis Proposal}
\author{
        Mark Reuter
}

\documentclass[12pt]{article}
\usepackage{mathrsfs}
\usepackage[utf8]{inputenc}
\usepackage[english]{babel}
\usepackage{amsthm}
\usepackage{amssymb}
\usepackage{amsmath}
\usepackage{tabularx}
\usepackage{ltxtable}

\newtheorem{theorem}{Theorem}[section]
\newtheorem{lemma}[theorem]{Lemma}

\theoremstyle{definition}
\newtheorem{definition}[theorem]{Definition}

\theoremstyle{remark}
\newtheorem*{remark}{Remark}

\usepackage[letterpaper, margin=1in]{geometry}
\usepackage{setspace}
\singlespacing

\usepackage{mathtools}
\DeclarePairedDelimiter\ceil{\lceil}{\rceil}
\DeclarePairedDelimiter\floor{\lfloor}{\rfloor}

\begin{document}
\maketitle

\singlespacing
\section{Introduction}
\doublespacing

\paragraph{Goal}

Create an algorithm which produces an eroded heightmap at a good-enough quality in comparison to the traditional erosion algorithms while not having the associated performance penalties.

\paragraph{Problem}

Perlin noise and fractal Brownian motion do not create realistic looking mountains and valleys. Water and thermal erosion algorithms are used on heightmaps generated by those noise functions to produce more realistic heightmaps, but the erosion algorithms suffer from requiring potentially thousands of iterations limiting the real time possibilities of these algorithms.

\paragraph{Proposed Solution}

I propose the use of noise editing to produce pregenerated noise tiles with eroded qualities. This will be done by feeding samples from precomputed erosion heightmaps into the noise editing algorithm, then estimating a close fit with a limited number of gradient vectors. Then at run time, a noise height map will be produced based on these noise tiles, and a fewer number of iterations of the erosion algorithm will be (1986, 2002)applied.

Throughout this paper, I will use bold letters with or without subscripts to denote vectors, and 

\singlespacing
\section{Related Works}
\doublespacing

\subsection{Noise Functions}

\paragraph{Perlin Noise}

Perlin Noise\cite{perlin1985} has been an ubiquitous noise function in compter graphics since its unveiling in 1985. Originally this function was designed for use in texture synthesis\cite{perlin1985, perlin1989}, but its applications to other areas of procedural generation have been well explored\cite{lagaeetal2010, smelik2014}. Perlin noise is a lattice gradient noise function. A \textit{lattice gradient noise function} generates noise by interpolating vectors, called \textit{gradients}, defined at the integer lattice points. Perlin's original interpolation function used the cubic polynomial $3t^2-2t^3$\cite{perlin1985}, but this function is discontinuous in the second derivative at integer points. For this reason, Perlin changed the interpolation function to the quintic polynomial $interp(t)=6t^5-15t^4+10t^3$, and this function has a zero first and second derivative at both $t=0$ and $t=1$\cite{perlin2002}.

Intuitively, Perlin Noise is computed as the weighted sum of $2^N$ gradients where $N$ is the dimension of the input vector $\textbf{v}$. The standard method for computing Perlin Noise defines an indexed set $G$ of gradients and a permutation $P$ of $\{1, ..., H\}$. $H$ can be any arbitrary positive integer, but implementations of Perlin Noise generally set $H=256$. The standard method for producing the weighted sum of gradients uses the $2^N$ integer lattice points surrounding $\textbf{v}$, named $\textbf{z}_1$, ..., $\textbf{z}_{2^N}$, where each $\textbf{z}_i = \langle z^i_1, ..., z^i_N \rangle$, and the gradients at the integer lattice points are obtained by computing the following equations:

\begin{align*}
	\textbf{g}_i = G[(P[z^i_1 \; \textnormal{mod} \; H] \oplus ... \oplus P[z^i_N \; \textnormal{mod} \; H]) \; \textnormal{mod} \; |G|] \\
	\textnormal{for each } 1 \leq i \leq 2^N, \\
	\textnormal{where } \oplus \textnormal{ is the bitwise exclusive-or.}
\end{align*}

These gradients $g_i$ are the ones on which Perlin Noise computes a weighted sum. The weights for use in this weighted sum are obtained using the interpolation function applied on the distances, named $\textbf{d}_1$, ..., $\textbf{d}_{2^N}$, between the integer lattice points to the input vector $\textbf{v}$. Specifically,the weights are the vectors $\textbf{c}_i$, that are computed as follows:

\begin{align*}
	\textbf{d}_i &= \textbf{v} - \textbf{z}_i\\
	\textbf{c}_i &= \textbf{d}_i * \prod_{j=1}^{N}(1-interp(|d^i_j|))
\end{align*}

Finally, Perlin Noise is computed as the sum of the dot product between the weight vectors and gradients:

\begin{equation*}
	perlin(\textbf{v}) = \sum_{i=1}^{2^N}(\textbf{c}_i \cdot \textbf{g}_i)
\end{equation*}

Perlin Noise has a regular zero-crossing property; if $\textbf{v}$ is identical to one of the integer lattice points, then $perlin(\textbf{v})=0$. The integer lattice points are not the only points at which Perlin Noise evaluates to zero; depending on definitions of $G$ and $P$, Perlin Noise can evaluate to zero elsewhere. Regardless of $G$ and $P$, though Perlin Noise always evaluate to zero at integer lattice points.

A noise function is \textit{band-limited} if it only produces frequencies within a band. Perlin Noise is not fully band-limited\cite{cookderose2005}. Patterns are constructed from \textit{bands} of noise which are limited to specific frequency ranges. Cook and DeRose have demonstrated that Perlin Noise produces both high and low frequencies outside of a band, and this leads to aliasing problems. (**elaborate**)

The choice of the gradients in the set $G$ can impact the quality of Perlin Noise in respect to the range of expressed noise values and band-limitedness. Originally,  Yoon et al.\cite{yoonetal2004} demonstrated that a 1-dimensional gradient set of two-thirds positive values and one-thirds negative values does not produce noise in the full range of $[-1, \; 1]$. Later, Yoon and Lee\cite{yoonlee2008} demonstrated that the choice of gradient set can produce more or less band-limited Perlin Noise.

\paragraph{Simplex Noise}

Ken Perlin\cite{perlin2001,gustavson2005} introduced another lattice gradient noise function called Simplex Noise. It is an improvement to his previous noise generation algorithm. Simplex Noise deviates from Perlin Noise in two key aspects. First, the integer lattice is replaced by the simplex lattice. Second, radial attenuation is used instead of polynomial interpolation on the gradients. These two modifications lower the computational complexity of the noise function in higher dimensions and remove the directional artifacts of Perlin Noise.

A \textit{$N$-simplex} is the simplest shape that tessellates $N$-dimensional space, and the \textit{simplex lattice} is the arrangement of these shapes to fill $N$-dimensional space. For 1-dimensional space, the simplex lattice is no different than the integer lattice. The simplex lattice in 2-dimensional space is a tessellation of equilateral triangles, and for 3-dimensional space, the simplex lattice is a tesselation of slightly skewed tetrahedra.

Like Perlin Noise, Simplex Noise is the weighted sum of $N+1$ gradients where $N$ is the dimension of the input vector $\textbf{v}$. Likewise, the standard method for computing Simplex Noise defines an indexed set of gradients $G$ and a permutation $P$ of $\{1, ..., H\}$, and implementations of Simplex Noise typically set $H=256$.

(**write a paragraph on the skew function**)

Simplex Noise must decide which simplex the input vector lies(**rewrite**). Simplex defines a skew function that squishes the simplex lattice to align with the integer lattice (**Define skew function**). After the alignment, $N$ factorial simplexes occupy the same space as a $N$-hypercube (**rewrite**). The simplex containing the input vector $\textbf{v}$ is decided by the ordering of the components of $\textbf{v}$. It should be noted that when skewed, each vertex of a simplex aligns with an integer lattice point. This fact is used to obtain the $N+1$ gradients. Let $\textbf{p}_1, ..., \textbf{p}_{N+1}$ be integer lattice points on which the vertices of the skewed simplex containing $\textbf{v}$ lie. Then,

\begin{align*}
	\textbf{g}_i=G[(P[p^i_1  \; \textnormal{mod} \; H] \oplus ... \oplus P[p^i_{N}  \; \textnormal{mod} \; H]) \; \textnormal{mod} \; |G|] \\
	\textnormal{for each } 1 \leq i \leq N+1, \\
	\textnormal{where } \oplus \textnormal{ is the bitwise exclusive-or.}
\end{align*}

Further, Simplex Noise weights the contribution from the gradients using radial attenuation. Conceptually, radial attenuation is like a drop of water falling in a still pool. A ripple expands out from the drop, and it lessens in magnitude the further it gets from the drop. Radial attenuation weights the contributions of the gradients in a similar manner. The contribution of a gradient is strongest the closer the noise is sampled to it, and the contribution falls off to zero before the input vector crosses the boundary into a simplex further away. Formally stated, (**check this**) with $\textbf{p'}_i$ the location of the $N+1$ vertices of the simplex containing the input vector $\textbf{v}$ and the two place function symbol, $e$, the Euclidean distance between two vectors squared, the weight vectors $\textbf{c}_i$ are defined as follows (**make sure to define the d vectors**):

\begin{align*}
d_i &= max\{0, 0.6-e(\textbf{p'}_i,\textbf{v})\} \\
\textbf{c}_i &= d_i^4 * (\textbf{p'}_i - \textbf{v})
\end{align*}

Finally, the simplex noise function is computed as the sum of the dot product between the weight vectors and gradients along with a constant multiplication $C$ to yield noise in the correct range. 

\begin{equation*}
	simplex(\textbf{v}) = C * \sum_{i=1}^{N+1}(\textbf{c}_i \cdot \textbf{g}_i)
\end{equation*}

Like Perlin Noise, Simplex Noise has a regular zero-crossing property; if the input vector $\textbf{v}$ is identical to one of the simplex lattice points, then $simplex(\textbf{v})=0$. Also, these are not necessarily the only points at which Simplex Noise evaluates to zero; depending on the definitions of $G$ and $P$, Simplex Noise can evaluate to zero elsewhere. Regardless of $G$ and $P$, though Simplex Noise always evaluates to zero at integer lattice points.

To my knowledge, the impact of the choice of the gradient set $G$ using the techniques of Yoon and Lee has not been done in respect to Simplex Noise. As well, study of the band-limitedness of Simplex Noise using the techniques of Cook and DeRose is not complete.

\paragraph{Fractal Sums}

(**return here**)

Fractal Sums are a method of producing fractal Brownian motion\cite{lagaeetal2010}. The Fractal Sum method computes the contribution of multiple octaves of noise. An \textit{octave} of noise is a pair $\langle f, a \rangle$ where $f$ is called the frequency and $a$ is called the amplitude. Implementations of Fractal Sums generally (**work here**) compute $F$ octaves of noise $\langle f_1, a_1 \rangle, ..., \langle f_F, a_F \rangle$, and each octave is related to the next by a lacunarity, $l$, and gain, $g$, as follows:

\begin{align*}
	f_1 &= 1, \\
	f_{i+1} &= f_i * l, \\
	a_1 &= 1, \\
	a_{i+1} &= a_i * g.
\end{align*}

In addition to these definitions of the octaves, typical implementations of Fractal Sums will define the lacunarity $l$ and gain $g$ in terms of each other. An implementation of Fractal Sums is called a $\frac{1}{l}$\textit{-noise} if $g=\frac{1}{l}$. Finally, the definition of Fractal Sums to be defined in terms of some one place noise function $noise$:

\begin{equation*}
	fractal(\textbf{v}) = \frac{\sum_{i=1}^{F}(a_i * noise(f_i * \textbf{v}))}{\sum_{i=1}^{F}a_i}.
\end{equation*}

This definition can have a regular zero-crossing property. For example, if the function $perlin$ as defined above is used with an integer lacunarity, then this function will cross zero at each integer lattice point. The regularity of this property can be modified by choosing a non-integer lacunarity.

\subsection{Noise Editing}

Yoon et al.\cite{yoonetal2004} presented a preprocessing algorithm called Noise Editing which modifies the gradient vectors used by Perlin Noise to evaluate a user specified constant at a given point. This algorithm is based on the observation that the Perlin Noise function can be rewritten as a weighted sum of $2^N$ gradient vectors where $N$ is the dimension of the noise. For any input vector $\textbf{v}$, the computation of the weight vectors in Perlin Noise is fixed by the algorithm. On the other hand, the gradients are chosen by definition. Yoon et al. noticed that a precisely chosen gradient set and permutation can force Perlin Noise to evaluate specified noise values for specified input vectors. This algorithm can be separated in two steps. First, a solution set is defined. Second, the optimal solution is chosen.

The first part of this algorithm is to define a solution set. Noise Editing defines a set of tuples $V = \{ (\textbf{v}_1, h_1), ..., (\textbf{v}_n, h_n) \}$ where each $\textbf{v}_i$ is a $N$-dimensional vector and each $h_i$ is a value in the range $-1$ to $1$. These pairs represent input vectors and the desired noise value for Perlin Noise to evaluate at that vector. The algorithm then defines a set of equations $E$ as follows:

\begin{equation*}
	E = \{ perlin(\textbf{v}) = (\textbf{c}_1 \cdot \textbf{g}_1) + ... + (\textbf{c}_{2^N} \cdot \textbf{g}_{2^N}) = h : (\textbf{v}, h) \in S \}
\end{equation*}

where each $\textbf{g}_i$ is a $N$-dimension gradient left variable and each $c_i$ is a weight vector as defined in the Perlin Noise section. Multiple input vectors can fall within the same cell of the integer lattice or within adjacent cells of the integer lattice. This implies that these input vectors consider the same or some overlapping gradients. The names $\textbf{g}_i$ in the equations must be chosen in the equations of $E$ to assure that the same gradient is considered for any reference to the same lattice point. Finally, $E$ forms a system of equations. Any solution to this system of equations will force Perlin Noise to evaluate the specified noise value $h_i$ for the given input vector $\textbf{v}_i$.

The second part of the algorithm concerns multiple solutions to the system of equations $E$. Some systems have multiple or even infinitely many solutions, and the choice of gradients can produce higher or lower quality noise. Yoon et al. define an optimization problem based on the $\chi^2$ test. The $\chi^2$ test subdivides a gradient set into $K$ buckets which all cover a uniform range of gradients. For a uniform random gradient set, each of these buckets will contain the same number of elements. But a solution chosen by Noise Editing does not guarantee a uniform distribution. For each $1 \leq k \leq K$, let $Z_k$ represent a bucket from the uniform random distribution and $Y_k$ represent a bucket from a chosen solution. Define the the value $D^2$ as follows:

\begin{equation*}
	D^2 = \sum_{k=1}^{K}(|Y_k| - |Z_k|)^2
\end{equation*}

Noise Editing computes the $D^2$ value for every solution or selection of solutions. Then, Noise Editing chooses the gradient set with the minimal $D^2$ value, or the gradient set closes to a uniform random distribution.

This algorithm can be generalized to several variations. Yoon et al. presented Noise Editing using Fractal Sums of Perlin Noise. As well, Yoon and Lee\cite{yoonlee2008} extended the optimization problem to include the weighted sum of the $x^2$, autocorrelation, and band-limitedness tests. Further, Cheng et al.\cite{cheng2014} argued that the Kolmogorov-Smirnov goodness of fit test provides better results than the $\chi^2$ test.

Noise Editing suffers from the zero-crossing property of Perlin Noise. An input vector which lies on an integer lattice point must evaluate to zero. If $\textbf{z}$ is an integer lattice point, no solution can exist for the pair $(z, h)$ where $h$ is non-zero. Yoon et al. gave two solutions to this problem. The first is noise field shifting. This solution shifts each vector in $V$ to avoid vectors which lie on integer lattice points. The second is noise field scaling. This solution applies a scalar multiplication to the noise values to fit the specifications of $V$. Neither of these are complete solutions as unsolvable cases still exist in both solutions.

\subsection{Tiling}

\paragraph{Wang Tiles}

Wang\cite{wang1961} presented a class of tiles, refered to as Wang tiles, which have colored edges and the stipulation that two tiles placed next to each other must have the same color on the adjoining edge. Wang tiles may not be rotated to fit. There are two basic questions concerning sets of Wang tiles. First, does the set fully tile Euclidean space? Second, does the set have a periodic placement of the tiles? Wang conjectured that if a set of Wang tiles filled Euclidean space, then there was a periodic ordering of the tiles. This conjecture was originally proved wrong in 1966 by Berger\cite{berger1966}, and Jeandel and Rao\cite{jeandelrao2015} found a set as small as 11 tiles using 4 colors which only has aperiodic tilings. They also proved that no set smaller than 10 tiles with 3 colors can force an aperiodic tiling of Euclidean space.

\paragraph{Corner Tiles}

Corner tiles\cite{lagaedutre2006,lagaeetal2006} were presented as an alternative to Wang tiles. Rather than coloring the edges, corner tiles color the corners with similar restrictions on which tiles can be placed adjacent to each other. Corner tiles have a distinct advantage to Wang tiles. Lagae and Dutr{\'e}\cite{lagaedutre2006} show that Wang tiles can cause discontinuities at the corners because diagonal tiles are unrelated to each other. Also corner tiles have much more efficient hashing algorithms similar to the hashing function used in standard Perlin Noise. Similar questions about sets of Wang tiles can be asked about sets of corner tiles. Corner tiles are less researched by the mathematical community, and as a result, the smallest set of aperiodic corner tiles known is 30 tiles using 6 colors found by Nurmi\cite{nurmi2016}.

\subsection{Erosion Algorithms}

Musgrave et al.\cite{musgraveetal1989} present water and thermal erosion algorithms for generating more realistic heightmaps. Both of these algorithms are iterative and can require thousands of iterations to produce realistic heightmaps which limits their use in the real-time generation of terrain\cite{smelik2014}. Olsen\cite{olsen2004} discusses several speed optimizations, but they reduce the quality of the produced heightmaps. Bene{\u s} and Forsbach\cite{benes2001} improved the realism of the origional algorithms by introducing material properties to the heightmaps, but these improvements come at a loss of performance. Gain et al.\cite{gain2015} implement the erosion algorithms with material properties and other geometric constraints on a GPU to greatly improve the performs of these algorithms. They produced 1024x1024 heightmaps in as little as 165ms. 

\subsection{Noise Tiles}

Maung\cite{maung2016} explored the use of Wang tiles for the procedural generation of textures as a means to reduce memory footprint of procedural textures. He introduces the concept of a Noise Tile to achieve these ends. \textit{Noise tiles} are $n \times n$ arrays of indices and a color vector $\langle n, e, s, w \rangle$ which associates a color with each edge of the tile. These tiles are defined with the same adjacency rules as a Wang tile. Maung's then modifies the standard method for computing Perlin Noise. In addition to the gradient set $G$ and permutation $P$, a set of Noise Tiles $N$ indexed by the color vector is defined. Then, the integer lattice is segmented into $n \times n$ regions, and a Noise Tiles is associated with each region. The permutation $P$ is used to decide the edge colors of each region. The only formal modifications Maung makes to the Perlin Noise function is in gradient selection. Let $T$ be the Noise Tile associated with the region in which the input vector $\textbf{v}$ lies. Likewise, let $\textbf{p}_1, ..., \textbf{p}_4$ be the integer lattice points surrounding $\textbf{v}$. Select each $\textbf{g}_i$ as follows:

\begin{align*}
	\textbf{g}_i = G[T[p^i_1][p^i_2]]
\end{align*}

The concept of a Noise Tile can be extended to generate Fractal Sums. A \textit{Fractal Noise Tile} is a sequence of $m$ noise tiles $T_1, ..., T_m$ where $T_1$ is a $n_1 \times n_1$ noise tile and each tile $T_{i+1}$ is a $(n_i * f) \times (n_i * f)$ noise tile where $f$ is the lacunarity of the fractal sum. The formal modifications follow the required modifications for standard noise tiles except that for each octave $i$ the noise tile $T_i$ is used.

There are limitations on noise tiles defined this way. Each index at a corner of a noise tile must be the same. Without this stipulation, there could exist discontinuities at the intersection of four regions because there is no correlation between two tiles which are diagonally adjacent.

\singlespacing
\section{Proposed Work}
\doublespacing

\paragraph{Simplex Noise Editing}

I propose a modification to Noise Editing which uses Simplex Noise instead of Perlin Noise. This modification requires the redefinition of the set of equations $E$. Given a set of tuples $V = \{ (\textbf{v}_1, h_1), ..., (\textbf{v}_n, h_n) \}$ where each $\textbf{v}_i$ is a $N$-dimensional vector and each $h_i$ is a value in the range $-1$ to $1$. $E$ can be defined in terms of Simplex Noise by consider $N+1$ weight vectors and gradients as follows:

\begin{equation*}
E = \{ simplex(\textbf{v}) = (\textbf{c}_1 \cdot \textbf{g}_1) + ... + (\textbf{c}_{N+1} \cdot \textbf{g}_{N+1}) = h : (\textbf{v}, h) \in S \}
\end{equation*}

where each $\textbf{g}_i$ is a $N$-dimension gradient left variable and each $\textbf{c}_i$ is a weight vector as defined in the Simplex Noise section. The rest of the algorithm is left unchanged.

\paragraph{Corner-Colored Noise Tiles}

I propose the use of Corner Tiles instead of Wang Tiles to define Noise Tiles. Rather than associating a color with each edge of the tile, I propose associating a color with each corner of the tile in accordance with the rules of Corner Tiles. This is to remove the stipulation that the index at each corner of a Noise Tile must be the same. Corner Tiles assure correlation between diagonally adjacent tiles, so the discontinuity problem in Noise Tiles is removed. Little modification to the definition of a Noise Tile is necessary, and I propose the generalization of this approach to Fractal Noise Tiles.

\paragraph{Simplex Noise Tiles}

I propose extending the concept of a Noise Tile to Simplex Noise. Simplex Noise in the second dimension uses a triangular lattice. This allows for Simplex Noise Tiles to use half the memory as a standard Noise Tile. The concept of a Simplex Noise Tile is a minor modification to standard Noise Tile. A Simplex Noise Tile is a $n \times n$ array of indices and a color vector $\langle c_1, c_2, c_3 \rangle$ which associates a color with each corner of a triangle. I introduce the concept of an \textit{upper} and \textit{lower} triangle. Simplex Noise skews the lattice to align with the integer lattice for the selection of gradients. This causes two right triangles to fill a cell of the integer lattice. I call the top of these two triangles an \textit{upper triangle} and the bottom a \textit{lower triangle}. The separation of upper and lower triangles allows for two triangles to share the same corner colors, and so this allows the upper and lower triangles with the same corner colors to fit into a single $n \times n$ Simples Noise Tile. Further, this allows a Simplex Noise Tile to be indexed using the same algorithm as a standard Noise Tile.

\subsection{Zero-Crossing Solutions}

\paragraph{Modified Perlin Noise} 

The generation of Perlin Noise only considers gradients selected from the $2^N$ integer lattice points surrounding the input vector. When the input vector falls on an integer lattice point, the spline based interpolation function causes the contribution from the other $2^N-1$ gradients to fall to zero. And the weight vector is based on the distance vector from the input vector to the integer lattice point. The contribution from the integer lattice point on which the input vector falls is also zero because the distance from the integer lattice point to the input vector is also zero. I propose a solution which considers more integer lattice points when generating noise and changes the interpolation function used.

This solution generates noised based on the $4^N$ integer lattice points surrounding the input vector. As well, radial attenuation is used instead of splines to interpolate the contribution from the gradients. This function is still defined as the sum of the dot product between weight vectors and gradients. The gradients $\textbf{g}_i$ are obtained in the same manner as Perlin Noise. The only difference being the number of gradients selected. Where $e$ is a two place function symbol denoting the Euclidean distance squared and $\textbf{z}_i$ are the surrounding integer lattice points, the weight vectors are redefined as follows:

\begin{align*}
d_i &= max\{0, 2 - e(\textbf{z}_i,\textbf{v})\} \\
\textbf{c}_i &= d_i^4 * (\textbf{z}_i - \textbf{v})
\end{align*}

And the proposed noise function is defined as follows:

\begin{equation*}
	noise_1(\textbf{v})=C*\sum_{i=0}^{4^N}(\textbf{c}_i \cdot \textbf{g}_i)
\end{equation*}

where $C$ is a constant multiple to obtain noise values in the range $-1$ to $1$. The zero-crossing property of Perlin Noise is caused by the contribution of all gradients to the noise value falling off to zero at integer lattice points. At integer lattice points, some of the gradients contributing to the noise value produced by $noise_1$ do necessarily fall off to zero. As in the definition of Perlin Noise, the contribution from the gradient at the integer lattice point on which the input vector falls is zero. But this noise function has a distinct advantage over Perlin Noise. At integer lattice points, not all gradients contributing to the noise value fall off to zero. So it is possible for the sum of the contributions to be non-zero. But the noise values produced at the integer lattice points tend to produce local minimums in the noise. This is due to the fact that the integer lattice point coinciding with the input vector has the largest contribution to the noise, but it necessarily falls to zero.

$noise_1$ calculates the contribution of many gradients whose weights have fallen off to zero. This is due to the fact that The contribution from the corners of the $2N$-hypercube intersect the innermost $N$-hypercube, but it does not intersect much of this hypercube. This leads to unnecessary computations. 

\paragraph{Modified Simplex Noise}

A downside to the proposed noise function $noise_1$ is its growth rate. It grows in $O(4^N)$ with the dimension $N$. Perlin Noise also has an exponential growth rate $O(2^N)$, but this growth rate is lessened by the fact that noise is generally evaluated in the second and third dimension and is rarely evaluated in dimensions higher than the fourth dimension. So while Perlin Noise does have exponential growth rate, it only computes the contribution from 4, 8, and 16 gradients in the second, third, and fourth dimensions respectively. So in the dimensions which Perlin Noise is most used, it is less computationally intensive than the growth rate might suggest. But the growth rate of $noise_1$ does pose a problem in these dimensions. $noise_1$ computes the contribution of 16, 64, and 256 gradients in the second, third, and fourth dimensions respectively.

One of the original reasons Ken Perlin introduced Simplex Noise was to reduce the growth rate of his original algorithm. To mirror this progress, I propose a second noise function the reduce the growth rate of $noise_1$ while still maintaining a reduced zero crossing property. This approach, like Simplex Noise, is based on the simplex lattice, and it maintains the use of radial attenuation as its interpolation function. There are two advantages to using the simplex lattice over the integer lattice. First, the contribution from fewer gradients is calculated. This noise function calculates the contribution from $(N+1)^2+(N+1)$ gradients rather than $4^N$ gradients. So the contribution from 12, 20, and 30 gradients are calculated in the second, third, and fourth dimensions respectively. Second, $noise_1$ calculates the contribution of many gradients which have fallen off to zero. This noise function calculates less of these zero contributions making the simplex lattice more favorable to the integer lattice.

The formal changes to Simplex Noise are minimal. The same skew function as simplex noise is used to convert the simplex lattice to an integer lattice for gradient selection. The only difference when selecting gradients is that, in addition to selecting gradients from the simplex contain the input vector, gradients are selected from the lattice points adjacent to those from the containing simplex. This amounts to the original $N+1$ gradients and an additional $(N+1)^2$ gradients. Let each $\textbf{p'}_i$ be one of these gradients, and let the two place function symbol $e$ be the Euclidean distance squared. The modification to the weight vectors is as follows:

\begin{align*}
d_i &= max\{0, 1.2 - e(\textbf{p'}_i,\textbf{v})\} \\
\textbf{c}_i &= d_i^4 * (\textbf{p'}_i - \textbf{v})
\end{align*}

This noise function is then defined as the sum of the dot products between these weight vectors and the selected gradients $\textbf{g}_i$

\begin{equation*}
noise_2(\textbf{v})=C*\sum_{i=0}^{(N+1)^2+(N+1)}(\textbf{c}_i \cdot \textbf{g}_i)
\end{equation*}

\paragraph{Generation of Noise Edited Tiles}

The culmination of this project is the definition of a Noise Edited Tile. Common applications in procedural terrain generation use highly iterative procedures to create world geometry. This limits the real-time capabilities of these algorithms. I propose a construction of Noise Tiles to improve the performance of constructing world geometry. This solution is a preprocessing algorithm. First, a world geometry is generated as a basis for Noise Editing. Then regions of the geometry are selected for the construction of Noise Tiles. These regions are sampled at regular intervals for the use in Noise Editing to select gradients. Finally, the Noise Tiles are constructed based on the optimal solution found by Noise Editing.

\singlespacing
\section{Work Schedule}
\doublespacing

\LTXtable{\textwidth}{schedule.tex}

\bibliographystyle{plain}
\bibliography{proposal}

\end{document}